\subsection{Part II design: memory, capacity or a clock?}
\label{sec:methods-part2}

Part~I compares an LSTM with an MLP that is 3.9 times smaller and cannot tell how far into the
episode it is. Part~II separates those differences with controls that change one thing at a time
(\cref{tab:arms}). All share the encoder of \cref{sec:arch} and the training and evaluation of
\cref{sec:training,sec:eval}, with a budget of 300,000 interactions.

\begin{itemize}
  \item \textbf{MLP-wide + clock} (11,399 parameters) replaces the LSTM cell with a memory-free
  MLP of the same size, $33\to128\to32$, whose extra input is the time inside the episode,
  $t/32$ ($t=0$ at the first step). Memory-free arms get the clock because, under a time limit,
  an agent that can count steps has information a memory-free agent lacks
  \citep{pardo2018timelimits}.
  \item \textbf{LSTM-reset} (11,367) is the LSTM with its state $(h,c)$ set to zero before every
  step, in training and evaluation: the same parameters and optimiser, no memory. Its 4,096
  recurrent weights receive no signal, so it controls for recurrence, not for capacity.
  \item \textbf{stack4 + clock} (11,381) sees the last four observations concatenated (zeros
  before the fourth step) and the clock: a short, fixed memory without learned recurrence
  \citep{hausknecht2015drqn}. Its body is $33\to55\to32$.
  \item \textbf{MLP-wide, no clock} (11,401) is MLP-wide without the clock input
  ($32\to130\to32$). It separates width from the clock. This arm was added after an interim look
  at 46 of the 80 phase-1 runs and is analysed as exploratory.
  \item \textbf{LSTM} and the small \textbf{MLP} are those of Part~I.
  \item Two \textbf{mechanism arms}, added after the confirmatory result and exploratory, know
  the time but have no memory: \textbf{LSTM-reset + clock} (11,495 parameters) feeds the clock to
  the cell of LSTM-reset, and \textbf{MLP + clock} (2,951) feeds it to the small MLP's encoder.
  They run on the same first 10 seeds with the protocol's learning rate, and their checkpoints
  were played on the same sealed panel under a separate evaluation record.
\end{itemize}

\paragraph{Learning rates.}
The new memory-free arms might need a different learning rate from the protocol's $10^{-3}$.
Before phase~1, MLP-wide + clock and stack4 + clock were each trained with $10^{-3}$,
$3\cdot10^{-4}$ and $3\cdot10^{-3}$ on three tuning seeds that are excluded from every
evaluation \citep{eimer2023hyperparameters}. A rule fixed in the plan chose the rate with the
highest mean validation AUC (ties to the protocol's rate) and wrote the choice to the chain's
journal before phase~1 started: $10^{-3}$ for MLP-wide (AUC 0.62 versus 0.58 and 0.12) and
$3\cdot10^{-4}$ for stack4 (0.42 versus 0.31 and 0.12). LSTM, LSTM-reset and the small MLP keep the
protocol's rate; MLP-wide without clock uses MLP-wide's.

\paragraph{Seeds.}
Twenty new training seeds come from \texttt{SeedSequence(20261020)}, excluding the pilot seeds,
the 30 seeds of Part~I and the tuning seeds. LSTM, MLP-wide + clock and LSTM-reset run on all 20;
the small MLP, stack4 + clock and MLP-wide without clock run on the first 10, so every arm is
paired with every other on the seeds they share. Within each seed, the order of the arms is a
fixed permutation drawn from the master seed. Every run is re-evaluated three times as in
\cref{sec:eval}. Phase~1 has 80 runs and the clock control 10.

\paragraph{A sealed panel of new maps.}
The confirmatory score is measured on 64 task-B maps that no run has seen and no study panel
contains. They were built without loading any model, by walking map seeds from 3,150,000 and
keeping maps whose physical identity differs from all 192 distinct maps of the study's panels.
For each run we play its validation-selected checkpoint (\texttt{best\_B}, chosen over the 300,000
interactions) on every sealed map with four new sampling seeds, 256 episodes per run. The panel
was opened once, after the analysis below was committed; the opening is recorded with the hash
of that analysis, and the tooling refuses to open it again with a different one.

\paragraph{Hypotheses and tests.}
The C01 design fixed, before any Part~II run, three principal contrasts and their thresholds:
\textbf{P1} (capacity), LSTM minus MLP-wide + clock is at least 15\pp; \textbf{P2} (recurrence),
LSTM minus LSTM-reset is at least 15\pp; and \textbf{P3}, a fully observable twin of the task,
which we did not run but still count in the family. The secondary contrast \textbf{H3} asks
whether a short fixed memory falls short of the two strong arms, that is, whether stack4 stays
below their average: $\text{LSTM} + \text{MLP-wide} - 2\,\text{stack4} > 0$. The unit of analysis is the training seed; maps are not replicates. For
each contrast we compute its value per seed on the seeds its arms share, an exact two-sided
sign-flip permutation test over all $2^n$ sign patterns, a BCa bootstrap 95\% interval with
10,000 resamples \citep{efron1993bootstrap} and the $t$ interval as a sensitivity check. Holm's
correction \citep{holm1979simple} uses $m=3$ for the principal family and $m=5$ for the secondary
one. A contrast is \emph{supported} when its Holm-adjusted $p<0.05$, its mean reaches the
threshold and the lower end of its interval is above zero. Contrasts added after the interim
look form an exploratory family that is reported without correction or verdict.

\paragraph{A probe for time.}
To see whether the LSTM's state carries the time step, we play each run's validation-selected
checkpoint on the 32 validation maps with two new sampling seeds and record, at every step $t$,
the LSTM's output after reading that step's observation. A ridge regression ($\lambda=1$,
standardised features) predicts $t$ from it, with four cross-validation folds split by map so that
no map appears in both training and testing. We report the out-of-fold $R^2$ and three
references: the same network at its initial weights, the raw observation (what is readable
without memory), and the trained state with shuffled labels \citep{hewitt2019control}.

The same analysis was first run on the development test panel of Part~I (the end-of-training
model, averaged over its four evaluations); those results were seen before the confirmatory
analysis was frozen. The analysis code is \texttt{ucl study-contrasts}, checked on synthetic data
with known answers.
